\section{Three-input rotations and the divisors they retain}
\label{sec:three-leaf-rotations}

The two-chart comparison holds at a fixed separation.  Changing a
bracketing changes which separations label the interval factors.
Three inputs are enough to determine the necessary localization and
to distinguish a coefficient equivalence from its action on states.

\subsection{The divisors of two three-leaf trees}

Let \(k\) be a field.  For three labelled positions set
\[
 t=z_1-z_2,\qquad x=z_2-z_3,\qquad u=z_1-z_3=t+x,
 \qquad R=k[t,x].
\]
The bracketing \(((1,2),3)\), retaining position 2 at its first
merger, uses the differences \(t,x\).  The bracketing
\((1,(2,3))\), retaining position 3, uses \(x,u\).
A rotation here means this change of bracketing with the spatial
functions \(t,x,u\) fixed.  It does not mean renaming independent
coordinates without their relation \(u=t+x\).

For a nonzero \(d\in R\), write \(A_d=\mathcal I_R(d)\)
for the interval algebra of \eqref{sb:eq:general-interval}.
Its endpoint idempotents are \(e_0^d,e_1^d\), its overlap generator
is \(\theta_d\), and its augmentation is denoted by \(p_d\).
All tensor products use the cohomological Koszul sign.
The two tree complexes are
\[
 E_L=A_t\otimes_R A_x,\qquad
 E_R=A_x\otimes_R A_u.
\]
Their tensor factors record different interval coordinates; the
rotation must specify a map between these actual complexes.

\begin{proposition}[Failure over the original spatial ring]
\label{tr:prop:global-obstruction}
Over every field, \(E_L\) and \(E_R\) are not
quasi-isomorphic as complexes of \(R\)-modules.
\end{proposition}

\begin{proof}
Localize \(R\) at the prime ideal \((t)\).  In the local ring
\(R_{(t)}\), both \(x\) and \(u=t+x\) are units.
When \(d\) is a unit in a coefficient ring, the maps
\[
 i_d(a)=(a,a),\qquad p_d(a_0,a_1)=a_0,\qquad
 h_d(\ell\theta_d)=(0,\ell),\quad h_d(A_d^0)=0
\]
satisfy \(p_di_d=\id\) and
\(d_Ah_d+h_dd_A=\id-i_dp_d\).  Thus the localized right
tree retracts onto \(R_{(t)}\).  The localized left tree
retracts onto \(A_t\otimes_R R_{(t)}\), whose first
cohomology is
\[
 R_{(t)}[t^{-1}]/R_{(t)}\ne0.
\]
The class of \(\theta_t/t\) is nonzero because \(t^{-1}\)
is not in \(R_{(t)}\).  Localization preserves cohomology:
an equation or a boundary in a localization can be checked after
multiplying by one allowed denominator.  A quasi-isomorphism over
\(R\) would therefore give a quasi-isomorphism after this
localization, contradicting the different first cohomology groups.
\end{proof}

The same argument at \((u)\) interchanges the two trees.
Consequently a comparison on a common open subset cannot retain
either generic divisor \(t=0\) or \(u=0\).  This obstruction
occurs before any state algebra, tree homotopy or completion is added.

\subsection{Retaining the common collision}

Invert precisely the two differing divisors and put
\[
 S=R[t^{-1},u^{-1}],\qquad Q=S[x^{-1}].
\]
Every \(A_d\) in this section is extended from \(R\) to
\(S\).  The coefficient of \(\theta_x\) lies in \(Q\),
while its endpoint coefficients lie in \(S\).  In particular
\(\theta_x/x\) still represents a nonzero class in
\(H^1(A_x)=Q/S\).  Both \(A_t\) and \(A_u\) have all
coefficients in \(S\).

Let \(D\) be any differential graded \(S\)-algebra, with
ordered, possibly noncommutative multiplication.  Set
\[
 E_L(D)=A_t\otimes_S A_x\otimes_S D,
 \qquad E_R(D)=A_x\otimes_S A_u\otimes_S D,
 \qquad E(D)=A_x\otimes_S D.
\]
Let \(p_L,p_R\) evaluate the \(t\)- and \(u\)-intervals
at their first endpoints.  Let \(i_L,i_R\) insert the unit
of the removed factor.  They are strict algebra chain maps.
Define
\begin{equation}\label{tr:eq:rotation-maps}
 r=i_Rp_L:E_L(D)\longrightarrow E_R(D),\qquad
 l=i_Lp_R:E_R(D)\longrightarrow E_L(D).
\end{equation}

\begin{proposition}[Explicit coefficient retractions]
\label{tr:prop:retractions}
The maps in \eqref{tr:eq:rotation-maps} are chain homotopy
equivalences over every field.  They are natural in strict dg
algebra maps of \(D\), with no flatness assumption.  With
\(F_L=i_Lp_L\), \(F_R=i_Rp_R\), their homotopies are
\[
 h_L=h_t\otimes\id\otimes\id,
 \qquad
 h_R(a\otimes b\otimes v)=(-1)^{|a|}a\otimes h_u(b)\otimes v.
\]
They satisfy
\begin{equation}\label{tr:eq:coefficient-homotopy}
 dh_T+h_Td=\id-F_T,\qquad h_T^2=h_Ti_T=p_Th_T=0,
 \qquad T=L,R.
\end{equation}
Moreover, for homogeneous \(a,b\in E_T(D)\),
\begin{equation}\label{tr:eq:twisted-derivation}
 h_T(ab)=h_T(a)b+(-1)^{|a|}F_T(a)h_T(b).
\end{equation}
\end{proposition}

\begin{proof}
For one invertible interval, differentiation on degree zero gives
\(h_dd_A(a_0,a_1)=(0,a_1-a_0)\); on degree one,
\(d_Ah_d(\ell\theta_d)=\ell\theta_d\).  These are
\(\id-i_dp_d\).  The side identities follow directly.

For \eqref{tr:eq:twisted-derivation} in one factor, two degree-zero
inputs give zero.  A degree-zero element \(a=(a_0,a_1)\)
and a one-form \(b=\ell\theta_d\) give
\(h_d(ab)=(0,a_0\ell)=i_dp_d(a)h_d(b)\).
In the reverse order the result is
\((0,\ell b_1)=h_d(\ell\theta_d)b\).
For two one-forms both sides vanish, because
\((0,\ell)m\theta_d=0\).  These cases prove the identity.

Tensoring with a complex preserves the displayed chain homotopy
when the graded tensor signs are included.  In the derivation
identity, moving a state across the degree-minus-one homotopy gives
exactly the change in its tensor sign.  Thus the first-factor
identity tensors with any dg algebra.  A graded permutation moves
the \(u\)-factor to the first position and back, giving the
stated sign for \(h_R\).  These arguments use explicit maps,
not exactness of tensoring.  Finally \(lr=F_L\) and
\(rl=F_R\), proving the assertion.
\end{proof}

The homotopy sends \(\theta_t/t^N\) to
\(e_1^t/t^N\).  Its existence uses the changed endpoint ring
\(S\); this primitive was not available over \(R\).
No pole class on the excluded divisor is preserved by that change.
The remaining \(x\)-collision is unchanged.

\subsection{Three polynomial state inputs}

An actual binary state operation determines more than its scalar
denominators.  Let \(V\) be the one-species free state algebra
\eqref{fs:eq:state-algebra}, with the Hasse translation
\eqref{fs:eq:hasse-translation}.  In this section \(S\) denotes the
localized spatial ring above, while \(V\) denotes states.
Use the kernels \(K_{ab}\) and put
\(P_{ab}^{ij}=P_{ij}^{ab}\), with the left odd derivatives of
\eqref{fs:eq:pair-operators}.  Every sum is finite before translation.
Write
\(E_d(a\otimes b)=(\tau_da)b\).  The binary overlap operation is
\begin{equation}\label{tr:eq:binary-state}
 Y_d=E_d\left(1+\theta_d\sum_{a,b\ge0}K_{ab}(d)P_{ab}^{12}\right).
\end{equation}
The uncontracted coefficient means the same value at both endpoints.
Coefficients precede states, and products of coefficient operators
use the graded tensor rule.  Thus an odd \(P\) crossing a
\(\theta\) contributes a minus sign.
Both terms in \eqref{tr:eq:binary-state} are closed and of total
degree zero.  The output is taken in the product of state weights,
with a common finite spatial pole bound for each polynomial input.

Set \(z_3=0,z_2=x,z_1=u=t+x\).  The left iteration first
uses \(Y_t\) on inputs 1 and 2, then the \(x\)-interval.
The right iteration first uses \(Y_x\) on inputs 2 and 3,
then the \(u\)-interval.  They are initially expanded in
different rings:
\[
 \iota_L:Q\longrightarrow k[x,x^{-1}]((t)),
 \qquad
 \iota_R:Q\longrightarrow k[u,u^{-1}]((x)),\quad t=u-x.
\]
The first expands \(u^{-1}=(x+t)^{-1}\) in \(t/x\),
and the second expands \(t^{-1}=(u-x)^{-1}\) in \(x/u\).
Both maps are injective, because their restrictions to the
polynomial ring are injective and the denominators map to units.
No identification of the two Laurent targets is made.

\begin{example}[One odd input]
\label{tr:ex:three-states}
Take \(a=x_0,b=\eta_0,c=x_0\), and write
\(E=(\tau_u x_0)(\tau_x\eta_0)x_0\).
The first binary operation on the left is
\(Y_t(x_0,\eta_0)=(\tau_tx_0)\eta_0+\theta_t/t\).
The remaining contraction pairs \(\eta_0\) with the last
\(x_0\).  Hence its rational representative is
\[
 Y_L(a,b,c)=E+\frac{\theta_t}{t}x_0
                +\frac{\theta_x}{x}\tau_ux_0.
\]
On the right the first contraction gives \(\theta_x/x\).
Contracting the first \(x_0\) with \(\tau_x\eta_0\)
gives \(\sum_{j\ge0}x^ju^{-j-1}=t^{-1}\) in its
specified right Laurent domain.  Thus
\[
 Y_R(a,b,c)=E+\frac{\theta_u}{t}x_0
                +\frac{\theta_x}{x}\tau_ux_0.
\]
Over \(S\), the coefficient rotation satisfies
\[
 rY_L(a,b,c)-Y_R(a,b,c)
 =-\frac{\theta_u}{t}x_0
 =d\left(-\frac{e_1^u}{t}x_0\right).
\]
The primitive has exactly the same pole order in \(t\).
For the inputs \(x_0,x_0,\eta_0\), even the projected
outputs have a difference:
\[
 p_LY_L-p_RY_R
 =\frac{\theta_x}{u}\tau_xx_0
 =d\left(\frac{e_1^x}{u}\tau_xx_0\right).
\]
The coefficient \(u^{-1}\) is regular on the retained
\(x\)-divisor over \(S\), so this second primitive is legal.
\end{example}

For general polynomial inputs define three odd operators
\[
 D_{12}=\sum_{a,b}K_{ab}(t)P_{ab}^{12},\qquad
 D_{13}=\sum_{a,b}K_{ab}(u)P_{ab}^{13},\qquad
 D_{23}=\sum_{a,b}K_{ab}(x)P_{ab}^{23},
\]
and the translated multiplication
\(E(a\otimes b\otimes c)=(\tau_ua)(\tau_xb)c\).
In expressions such as \(\theta_x ED_{13}\), the coefficient
is placed first and the displayed state map acts on the three inputs.

\begin{theorem}[A direct three-input rotation homotopy]
\label{tr:thm:state-rotation}
Over every field, the two binary iterations have the rational
representatives
\begin{equation}\label{tr:eq:tree-state-operations}
 \begin{aligned}
 Y_L&=E\bigl(1+\theta_x(D_{13}+D_{23})\bigr)
             \bigl(1+\theta_tD_{12}\bigr),\\
 Y_R&=E\bigl(1+\theta_u(D_{12}+D_{13})\bigr)
             \bigl(1+\theta_xD_{23}\bigr).
 \end{aligned}
\end{equation}
They are degree-zero chain maps from \(V^{\otimes3}\) to the
weightwise outputs of \(E_L(V)\) and \(E_R(V)\) over \(S\).
For the rotation \(r\), a degree-minus-one homotopy is
\begin{equation}\label{tr:eq:state-rotation-homotopy}
 \begin{aligned}
 H={}&e_1^x ED_{13}
      -e_1^u E(D_{12}+D_{13})\\
     &+\theta_x e_1^u E(D_{12}+D_{13})D_{23},
 \qquad dH+Hd_{V^{\otimes3}}=rY_L-Y_R.
 \end{aligned}
\end{equation}
Every term of \(H\) uses a spatial coefficient already occurring
in the displayed operations.  It introduces no new pole or higher
pole exponent.  The assertion concerns polynomial inputs and their
rational representatives, not arbitrary completed input states.
\end{theorem}

\begin{proof}
The derivative of a product is the sum of derivatives on its factors,
with the left-odd-derivative signs.  Thus a contraction of a merged
left state with input 3 splits into the 13 and 23 contractions.
Translation changes the first separation from \(x\) to \(x+t\).
For its coefficients the required identity is
\[
 \sum_{m\ge0}\binom{a+m}{a}t^m K_{a+m,b}(x)
 =K_{ab}(x+t).
\]
Indeed,
\(\binom{a+m}{a}\binom{a+m+b}{a+m}
=\binom{a+b}{a}\binom{a+b+m}{m}\), and the integer
geometric-series formula for a negative power gives the equality.
For the right iteration the corresponding identity is
\[
 \sum_{m\ge0}\binom{b+m}{b}x^m K_{a,b+m}(u)
 =K_{ab}(u-x).
\]
These are formal identities in the two specified Laurent domains.
The translation law gives \(\tau_x\tau_t=\tau_u\).
The product rule, these two identities and the original input order
therefore give \eqref{tr:eq:tree-state-operations}, including the
double-contraction terms.  All identities hold over the integers
before reduction to \(k\).

First-endpoint projection kills every term containing the removed
interval form.  Hence
\[
 p_LY_L-p_RY_R=\theta_x ED_{13}.
\]
All coefficients of \(D_{13}\) lie in \(S\), so
\(H_0=e_1^x ED_{13}\) is defined in \(A_x\otimes V\)
and \(dH_0=\theta_x ED_{13}\).  Each operation is closed,
since its coefficients are diagonal constants or products of closed
interval forms.  Proposition~\ref{tr:prop:retractions} gives
\[
 rY_L-Y_R=i_RdH_0-(\id-i_Rp_R)Y_R
           =d(i_RH_0-h_RY_R).
\]
To compute the last expression with the factors ordered \(x,u\),
move the odd contraction past the preceding interval form:
\[
 Y_R=E+\theta_x ED_{23}+\theta_u E(D_{12}+D_{13})
       +\theta_x\theta_u E(D_{12}+D_{13})D_{23}.
\]
Here the operator-crossing minus and the exchange of the two forms
cancel.  Now \(h_R(\theta_u)=e_1^u\) and
\(h_R(\theta_x\theta_u)=-\theta_xe_1^u\).
This gives exactly \eqref{tr:eq:state-rotation-homotopy}.
Equivalently, differentiating its last term uses
\(d(\theta_xe_1^u)=-\theta_x\theta_u\), which verifies
the signed identity directly.

There are at most two contractions.  Every contraction differentiates
one of the original polynomial inputs, so only finitely many indices
occur before translation.  Each fixed output weight receives
finitely many translation terms.  The denominators are products of
the displayed \(K_{ab}\).  The primitive only replaces interval
forms by endpoint idempotents, proving the pole assertion.
\end{proof}

This state result uses the specified free algebra \(V\) and its
binary operation.  The coefficient retractions of
Proposition~\ref{tr:prop:retractions} allow arbitrary noncommutative
dg state algebras; they do not prescribe a collision operation on
an arbitrary such algebra.

\subsection{Finite bar words and their homotopies}

Suppose first that \(D\) is augmented over \(S\).  Give
\(E_T(D)\) the augmentation obtained by first-endpoint evaluation
and the augmentation of \(D\).  For any augmented dg algebra
\(E\), write \(\bar E\) for its augmentation ideal and put
\[
 B(E)=\bigoplus_{n\ge0}(s\bar E)^{\otimes_S n},
 \qquad |[a_1|\cdots|a_n]|=\sum_i|a_i|-n.
\]
The symbol \(s\) decreases total cohomological degree by one.
Words are finite.  The differential \(b=b_1+b_2\) is given by
\[
 b_1[a]=[da],\qquad b_2[a|c]=(-1)^{|a|+1}[ac],
\]
with prefix sign \((-1)^{\sum_{i<j}(|a_i|-1)}\) when
the local operator starts in position \(j\).
The Leibniz identity cancels the mixed terms in \(b^2\).
On three entries the two remaining terms are
\((-1)^{|c|}[(ac)e]\) and \(-(-1)^{|c|}[a(ce)]\).
They cancel by associativity.  Disjoint local operations cancel by
their degree-one signs, so \(b^2=0\).

Deconcatenation cuts every word into consecutive subwords:
\[
 \Delta[a_1|\cdots|a_n]
 =\sum_{j=0}^n[a_1|\cdots|a_j]\otimes[a_{j+1}|\cdots|a_n].
\]
Two cuts choose three consecutive subwords, proving coassociativity.
A differential or multiplication away from the cut belongs to one
tensor factor with its original prefix sign.  A multiplication across
a proposed cut removes that cut on both sides.  Thus \(\Delta\)
is a chain map.

\begin{theorem}[Rotation on algebraic bars]
\label{tr:thm:bar-rotation}
Over every field, \(B(r)\) and \(B(l)\) are mutually inverse
chain homotopy equivalences of the finite-word bars.  They are strict
coalgebra maps.  For \(T=L,R\), define
\begin{equation}\label{tr:eq:bar-homotopy}
 K_T[a_1|\cdots|a_n]
 =\sum_{j=1}^n(-1)^{\sum_{i<j}(|a_i|-1)}
 [F_Ta_1|\cdots|F_Ta_{j-1}|h_Ta_j|a_{j+1}|\cdots|a_n].
\end{equation}
Then
\[
 bK_T+K_Tb=\id-B(F_T),\qquad
 \Delta K_T=(B(F_T)\otimes K_T+K_T\otimes\id)\Delta.
\]
The homotopies have degree minus one and preserve word length.
\end{theorem}

\begin{proof}
The coefficient homotopy preserves the augmentation ideal because
its first endpoint is zero.  The maps in
\eqref{tr:eq:bar-homotopy} are therefore defined on reduced bars.
For \(b_1K_T+K_Tb_1\), terms on different entries cancel by
the tensor sign rule.  At the entry where \(h_T\) acts, the coefficient
identity gives \(a_j-F_Ta_j\).  The sum telescopes from the
word with no \(F_T\) factors to the word with all such factors.
Its value is \(\id-B(F_T)\).

For the multiplication terms, positions away from that entry
cancel.  On two consecutive entries the remaining expression is
\[
 (-1)^{|a|}[h_T(a)b]+[F_T(a)h_T(b)]
       +(-1)^{|a|+1}[h_T(ab)],
\]
which is zero by \eqref{tr:eq:twisted-derivation}.  The same
prefix sign multiplies this identity inside any word.  This proves
the chain homotopy equation in every length.

For the cut identity, a cut before the entry where \(h_T\) acts contributes
\(B(F_T)\otimes K_T\).  A cut after it contributes
\(K_T\otimes\id\).  The degree-minus-one sign for the
second tensor factor is the prefix sign in
\eqref{tr:eq:bar-homotopy}.  These are all cuts.  Finally
\(B(l)B(r)=B(F_L)\) and \(B(r)B(l)=B(F_R)\), proving
the assertion.
\end{proof}

\subsubsection{Spatial pole bounds}

For a rational scalar \(f\), measure poles on the three distinct
prime divisors \(t,x,u\) of \(R\).  A tuple
\(N=(N_t,N_x,N_u)\) bounds its poles when
\[
 f\in t^{-N_t}x^{-N_x}u^{-N_u}R.
\]
Use the intersection of this space with the appropriate endpoint
or overlap coefficient module.  On a word, take the linear span of
tensors for which the sum of the input tuples is at most \(P\).
The span is necessary: sums can have poles on different divisors.

The coefficient maps \(p,i,h\) merely evaluate endpoint
components, insert a unit, or move a scalar from an overlap to an
endpoint module where it is allowed.  They do not change that scalar.
Consequently \(B(r),B(l),K_T\) preserve the total pole tuple.
Multiplication adds tuples, so the bar differential also preserves it.
The bound is attained by \([\theta_t/t^N]\), whose homotopy
is \([e_1^t/t^N]\).

For general state coefficients this statement uses a specified
regular state lattice: an \(R\)-dg algebra whose scalar extension
is \(D\), or an equivalent filtration preserved by its product
and differential.  Take images of the displayed scalar lattices
tensored with that state lattice.  This definition does not require
flatness or injectivity of scalar extension.  The maps respect the
resulting image filtration.  Without such a lattice, the chain
comparison still exists, but an intrinsic pole order of a state
coefficient has not been specified.

\subsubsection{Module ends}

Let \(M,N\) be a right and a left \(D\)-module concentrated
in degree zero, with dg actions.  An augmentation of \(D\) is
not required.  Let \(E_T(D)\) act through first-endpoint
evaluation of both interval factors and the stated \(D\)-actions.
The two-sided bar is
\[
 B(M,E_T(D),N)=\bigoplus_{n\ge0}
 M\otimes_S(sE_T(D))^{\otimes_S n}\otimes_S N.
\]
It has the internal differential and multiplication signs above.
The left outer face has positive sign; the right outer face has
sign \((-1)^{\sum_i(|a_i|-1)}\).  Nonzero end actions
require degree-zero entries.  Associativity and the module identities
prove the face identities and hence the square-zero differential.

Apply \(r,l\) to the middle entries and the identity to the
ends.  Formula~\eqref{tr:eq:bar-homotopy} applies to the middle
word.  The new end terms cancel because \(h_T\) has zero
first endpoint and \(F_T\) has the same endpoint as the identity.
Thus the same homotopy equivalences hold for these two-sided bars.
They preserve the actual module actions.

In particular, take the finite mode algebra with generators
\(q_m,p_m\), \(m\ge1\), and relations
\[
 [p_m,q_n]=\kappa m\delta_{mn},\qquad
 [p_m,p_n]=[q_m,q_n]=0.
\]
Its left module is \(W=k[h_0,h_1,\ldots]\), with
\(q_m=h_{m-1}\cdot\) and
\(p_m=\kappa m\partial_{h_{m-1}}\).  Give \(h_j\)
weight \(j+1\).  The restricted dual
\(M=\bigoplus_E W_E^\vee\) is a right module under
\((\lambda a)(v)=\lambda(av)\): every finite mode word
has a fixed weight shift, so it preserves finite weight support.
After extension to \(S\), these are admissible ends in the
two-sided comparison.  No character of the mode algebra is used;
when \(\kappa\ne0\), \([p_1,q_1]=\kappa1\) forbids one.
All entries remain finite mode words.

For a degree-zero state algebra \(D\), these module bars also
have balanced cuts through its regular module:
\[
 \lambda[a_1|\cdots|a_n]v\longmapsto
 \sum_{j=0}^n\lambda[a_1|\cdots|a_j]1
     \otimes_D1[a_{j+1}|\cdots|a_n]v.
\]
The target is
\(B(M,E_T(D),D)\otimes_D B(D,E_T(D),N)\), with the
interval factors acting on \(D\) through first-endpoint
evaluation.  At two neighboring cuts, the last face of the left
factor and the first face of the right factor cancel after
balancing over \(D\).  Their nonzero endpoint actions require
degree-zero entries, so their prefix signs are opposite.  All other
faces lie on one side of a cut and have the same tensor signs.
Thus the cut is a chain map.  The rotation maps commute with these
cuts because they preserve the endpoint action and the inserted
unit of \(D\).  The ends are not identified with scalar units.

\subsection{Polynomial forms and the characteristic-zero comparison}

For a divisor \(d\), the polynomial-form coefficient algebra is
\[
 C_d^0=\{f(s_d)\in S[d^{-1}][s_d]:f(0),f(1)\in S\},
 \qquad C_d^1=S[d^{-1}][s_d]ds_d,
 \qquad df=f'(s_d)ds_d.
\]
Products are products of forms, with \((ds_d)^2=0\).
The augmentation evaluates the first endpoint.  For \(d=t,u\)
this is the full polynomial de Rham algebra over \(S\).
Define form trees \(C_L(D)=C_t\otimes C_x\otimes D\)
and \(C_R(D)=C_x\otimes C_u\otimes D\).
All tensor products here are over \(S\).

Assume characteristic zero for the remainder of this section.
First-endpoint projection on an invertible interval has the
polynomial homotopy
\[
 h^{\Omega}(g(s)ds)=\int_0^s g(v)dv,
 \qquad h^{\Omega}(f)=0.
\]
It satisfies \(dh^{\Omega}+h^{\Omega}d=\id-ip\) and
the same side conditions as \eqref{tr:eq:coefficient-homotopy}.
This homotopy is not a twisted derivation in general.  Its extension
to full bars therefore needs a correction.

\begin{lemma}[A finite correction on form bars]
\label{tr:lem:form-bar-homotopy}
Let \(H_0\) be the tensor-word formula
\eqref{tr:eq:bar-homotopy} using \(h^{\Omega}\), and put
\(F=B(ip)\).  On a word of length \(n\), define
\begin{equation}\label{tr:eq:perturbed-homotopy}
 H=H_0\sum_{j=0}^{n-1}(-b_2H_0)^j.
\end{equation}
Then \(bH+Hb=\id-F\).  The endpoint projections and
unit inclusions induce mutually inverse chain homotopy equivalences
between the two form-tree bars through \(B(C_x\otimes D)\).
All maps are defined for arbitrary augmented dg state algebras,
and preserve the stated total pole filtrations.
\end{lemma}

\begin{proof}
The tensor contraction gives
\(b_1H_0+H_0b_1=1-F\), \(H_0^2=0\), and
\(FH_0=H_0F=0\).  Projections and inclusions are strict
algebra maps, so \(F\) commutes with \(b_2\).
Also \(H_0b_2F=0\) and \(Fb_2H_0=0\): use that
\(b_2\) commutes separately with the bar inclusion and
projection, which \(H_0\) kills on the corresponding side.

Write \(d=b_1\), \(\delta=b_2\) and \(Q_0=1+\delta H_0\).
The inverse of \(Q_0\) is a finite geometric series on each
word, because \(\delta\) decreases length and \(H_0\)
preserves it.  Thus \(H=H_0Q_0^{-1}\).
The identities \(\delta^2=0\),
\(d\delta+\delta d=0\) and
\(dH_0+H_0d=1-F\) give
\[
 dQ_0-Q_0d=-\delta(1-F),\qquad
 dQ_0^{-1}-Q_0^{-1}d=Q_0^{-1}\delta(1-F)Q_0^{-1}.
\]
Consequently
\[
 \begin{aligned}
 (d+\delta)H+H(d+\delta)
 &=1-F+H_0Q_0^{-1}\delta
              \bigl(1-(1-F)Q_0^{-1}\bigr)\\
 &=1-F+H_0Q_0^{-1}\delta(\delta H_0+F)Q_0^{-1}
 =1-F.
 \end{aligned}
\]
The final term is zero by \(\delta^2=0\) and
\(H_0\delta F=0\).  This proves the chain identity without
a convergence argument.  Tensoring the explicit contraction with
\(D\) includes its differential and needs no flatness assumption.
Every correction uses each original spatial coefficient once in a
product.  Integration changes only the interval coordinate, so
multiplication adds the original pole tuples and introduces no
new divisor.  The same argument with the two outer module faces
included in \(\delta\) gives the two-sided assertion for the
module ends specified above.
\end{proof}

The form-to-overlap comparison can also be made compatible with the
restricted rotation.  Its linear component on one interval is
\[
 I(f)=(f(0),f(1)),\qquad
 I(g(s)ds)=\left(\int_0^1g(s)ds\right)\theta.
\]
For a block of \(n>1\) entries, it is zero unless every
coefficient is a one-form.  On such a block it is
\[
 [g_1ds\otimes v_1|\cdots|g_nds\otimes v_n]
 \longmapsto
 [J_n(g_1,\ldots,g_n)\theta\otimes v_1\cdots v_n],
\]
where
\[
 J_n(g_1,\ldots,g_n)=
 \int_{0\le s_1\le\cdots\le s_n\le1}
                 g_1(s_1)\cdots g_n(s_n)\,ds_1\cdots ds_n.
\]
For monomials this is
\(\prod_{j=1}^n(r_1+\cdots+r_j+j)^{-1}\).
Sum over all consecutive partitions of a word, applying this map
to each block.  Denote the resulting coalgebra map by \(\mathcal I_d\).

For completeness, its chain identity reduces to one output entry.
For all-one-form blocks, coefficient differentials and products vanish.
For a block with one function \(f\), integration of \(f'\)
between its neighboring variables cancels the two adjacent products.
At the left end the surviving term uses \(f(0)\); at the
right end it uses \(f(1)\).  These are exactly the two endpoint
actions in \(A_d\).  Two functions reduce to endpoint
multiplicativity in length two; all remaining higher components
vanish.  With states, the common suspended prefix sign is
\((-1)^{\sum_{i<j}|v_i|}\).  The sign from moving the state
past the next form gives the multiplication sign, while a leading
form gives the minus in its state differential.  State products
stay ordered.  Cutting partitions commutes with partitioning cuts,
which proves the coalgebra identity.

The linear map \(I\) is a chain homotopy equivalence.  A
reverse map is linear interpolation on functions and
\(W(\ell\theta)=\ell ds\).  Its homotopy on one-forms is
\(\int_0^s g-s\int_0^1g\), with zero endpoints.
Direct differentiation gives \(IW=\id\) and
\(dh+hd=\id-WI\).  Tensoring this explicit contraction
and inducting on finite word length proves that \(\mathcal I_d\)
is a quasi-isomorphism, with no flatness condition.  The proof also
applies to genuine module ends: higher form blocks have zero first
endpoint, and singleton endpoint actions are unchanged.

Apply these comparisons in the order \(t\), then \(x\),
on the left, and \(x\), then \(u\), on the right.  Denote
the results by \(\mathcal I_L,\mathcal I_R\).
Write \(r_C\) for the form-tree rotation obtained by endpoint
projection and unit inclusion, and \(r_A=r\) for the overlap
rotation.

\begin{proposition}[Compatibility of the two presentations]
\label{tr:prop:form-square}
In characteristic zero, on finite bar words,
\[
 B(r_A)\mathcal I_L=\mathcal I_R B(r_C).
\]
Both maps preserve cuts and total spatial pole bounds.  The same
identity holds with the stated module ends.
\end{proposition}

\begin{proof}
Projecting an interval to its first endpoint kills every higher
block, since its coefficient is \(\theta\).  The remaining
singleton blocks are first-endpoint evaluation of the original
forms.  Conversely, inserting a constant interval factor gives no
higher component in that factor.  These two observations show
\[
 B(p_L^A)\mathcal I_L=\mathcal I_x B(p_L^C),\qquad
 \mathcal I_R B(i_R^C)=B(i_R^A)\mathcal I_x.
\]
They hold term by term on consecutive partitions.  Their composition
is the required identity.  Endpoint actions at module ends are
unchanged by the same calculation.  Products and integrals obey the
previously proved coefficient bounds.
\end{proof}

In characteristic \(p>0\), let
\(f=(s_x-s_x^p)/x\in\bar C_x^0\).  Its endpoints are
zero and \(df=ds_x/x\).  The prescribed linear comparison
would give \(I(f)=0\) and \(I(df)=\theta_x/x\ne0\).
Thus it fails the chain equation on the length-one word \([f]\).
A coaugmented coalgebra map sends primitive length-one words to
length-one words: any component of length at least two is detected
by its first proper cut.  Higher components cannot repair the
failure.  In addition, polynomial forms on an invertible interval
have nonconstant closed functions \(s^p\), so endpoint
projection is not a quasi-isomorphism there.  These obstructions
do not affect the all-field overlap homotopies above.

\subsection{The common formal domain and its limits}

The ring \(S\) retains the generic divisor \(x=0\) while
excluding \(t=0\) and \(u=0\).  A formal domain retaining
this same collision is obtained by
\[
 t=\lambda,\qquad x=\epsilon\lambda,\qquad
 u=(1+\epsilon)\lambda.
\]
These substitutions define
\[
 S\longrightarrow T=k[[\epsilon,\lambda]][\lambda^{-1}],
 \qquad
 Q\longrightarrow T[\epsilon^{-1}]
       =k[[\epsilon,\lambda]][(\epsilon\lambda)^{-1}].
\]
Both maps are injective on the rational rings: the polynomial
substitution is injective, and the inverted functions map to units.
For example,
\[
 t^{-a}x^{-b}u^{-c}
 =\lambda^{-a-b-c}\epsilon^{-b}(1+\epsilon)^{-c}.
\]
Its expansion has common finite lower bounds in both scales.
The retained overlap algebra has endpoints in \(T\) and
one-forms in \(T[\epsilon^{-1}]\theta_x\).
More precisely, a scalar with pole tuple \(N\) maps into
\[
 \epsilon^{-N_x}\lambda^{-(N_t+N_x+N_u)}k[[\epsilon,\lambda]].
\]
Polynomial numerators add only nonnegative scale powers.  The
monomial with numerator one attains both displayed bounds, since
\((1+\epsilon)^{-N_u}\) is a unit with constant term one.

For fixed polynomial inputs in
Theorem~\ref{tr:thm:state-rotation}, only finitely many contraction
coefficients occur before translation.  Each has the displayed
lower bounds.  Hasse translation adds nonnegative scale powers.
Each fixed state weight and scale coefficient receives finitely many
terms.  Thus the state identity, coefficient homotopies and finite
bar formulas extend coefficientwise to this specified domain.
An \(x\)-pole of order \(b\) becomes an \(\epsilon\)-pole
of order \(b\); powers of the invertible \(\lambda\) are
tracked separately.

This common expansion is reached from the rational representatives.
It is not a substitution on every element of the original left
Laurent completion.  For example,
\[
 \sum_{n\ge0}t^n x^{-n^2}
 \in k[x,x^{-1}][[t]]
\]
would become
\(\sum_{n\ge0}\epsilon^{-n^2}\lambda^{n-n^2}\),
with no lower bound in either scale.  It lies outside
\(T[\epsilon^{-1}]\).  The right and left Laurent targets
remain distinct.  Extending scalars from \(S\) to a ring
where \(x\) itself is a unit would also erase the retained
class \(\theta_x/x\); such an extension is a further change
of collision boundary.

There is a separate state-completion restriction.  In the product
of state weights, \(\eta_j\) tends to zero.  The scalar
component of \(Y_x(x_0,\eta_j)\) is
\(\theta_x/x^{j+1}\), which does not tend to zero in the
discrete scalar-state coefficient space.  The construction therefore
does not extend that binary operation continuously to all completed
input states.

The coefficient rotation and the free-state homotopy compare the
specified three-leaf operations after changing the endpoint rings.
They do not identify the two entire coefficient complexes over
\(R\), by Proposition~\ref{tr:prop:global-obstruction}.
Nor do they construct morphisms in a geometric factorization
category.  Such morphisms must preserve collision pullbacks,
factorization equivalences and their higher refinements.

\begin{problem}[Coherent rotations with collision boundaries]
Construct a common geometric chain carrier retaining the required
collision strata for several binary trees.  Specify maps from the
tree presentations and homotopies for their rotations.  For four
inputs, compare the five bracketings and prove the boundary identity
for the resulting pentagon of homotopies.  If additional states or
coefficient complexes are used to retain the divisors lost over
\(S\), give their differential and their maps to the stated
two-chart presentations.  Neither a composition of the three-leaf
maps nor equality of rational scalar expansions supplies these data.
\end{problem}
